\section{OpenAI 的组织文化与竞争格局}

\subsection{人才密度与组织架构}

翁家翌赞同 Sam Altman 的说法："在一个人才密度极高的小团队里，任何平庸的表现都是不能被容忍的。"人才密度高可以自发涌现出意想不到的东西。

\begin{knowledgebox}{OpenAI 的信息流通机制}
OpenAI 保持创新能力的关键是\textbf{信息流通通畅}。Sam Altman 有专门的研究助理帮他了解最新的内部研究进展；Greg Brockman 对整个 infra 底层几乎都参与过。领导层对技术细节的深入了解，确保了决策层和执行层的信息一致。
\end{knowledgebox}

从 280 人到 3000+ 人，OpenAI 面临的核心挑战是如何保持小团队的创新效率。翁家翌认为概率在下降但没有下降太厉害——关键是能划分出小团队专门做研究，同时简化组织架构、取消不合理的 meeting。

\subsection{Sam Altman 被开除事件的内部视角}

2023年11月 Sam Altman 被董事会开除的事件中，翁家翌提供了内部视角：底下干活的人\textbf{完全不知道发生了什么}，非常 surprise。董事会对员工缺乏透明度，决策过程不透明。核心原因是 Ilya 等董事会成员对 Sam 的\textbf{不信任}——不是不信任某个技术决策，而是不信任这个人。

翁家翌的核心诉求是\textbf{组织稳定}——"不要再出现一次公司差点倒闭的事"。组织架构的稳定有利于快速向前推进。在他看来，实现 AGI 的最大机遇和挑战都是一个词：\textbf{执行}——"对着正确的方向执行，只要能执行就好。"

\subsection{DeepSeek 的竞争与 OpenAI 的挑战}

翁家翌坦率地分析了 DeepSeek 的竞争优势：

\begin{warningbox}{大公司的规模劣势}
翁家翌认为 DeepSeek 的效率远高于 OpenAI——代码库小、沟通成本低、专注于特定 use case。而 OpenAI 要同时考虑很多 use case，各方面 trade off，组织大了必然面临效率下降。但 OpenAI 在用户反馈等其他维度有优势，这是 trade off。\textbf{每个公司都会变慢——看谁不那么慢。}
\end{warningbox}

翁家翌透露 DeepSeek 开源后，OpenAI 内部确实重新评估了开源策略。John Schulman 甚至问过他要不要把 RL Infra 开源，但出于公司利益考虑，翁家翌当时认为不太合适。不过他表示如果有无限资源，他"当然会很开心"地开源。

这里涉及到一个博弈论问题：OpenAI 如果开源最前沿的技术，其他公司会马上追上来，然后 OpenAI 可能融不到资、没有人持续输血。翁家翌认为"不是所有人都同条心的"——哪怕 OpenAI 想为了 AGI 造福全人类而开源，也有人只想借此赚钱。为了防着这种情况，OpenAI 不得不闭源。

\begin{knowledgebox}{开源的博弈论困境}
翁家翌指出了开源策略的核心矛盾：理论上存在一条路径——开源并接受社区反馈可以更好地实现 AGI。但实际执行非常困难，因为（1）你是第一名，开源了别人马上变第一；（2）别人再训练一下就超过你；（3）导致你融不到资。这是一个典型的囚徒困境——个体理性导致集体次优。
\end{knowledgebox}

\subsection{关于"OpenAI 的 Open"}

翁家翌对 OpenAI 的"Open"做了独特的解读：

\begin{importantbox}{OpenAI 的 Open：面向普通人}
OpenAI 的 Open 不是对其他大模型公司的 Open，而是\textbf{对普通人的 Open}——以尽可能便宜的价格（甚至免费 tier）让所有人都能接触到最先进的技术。"你丢一个裸的模型权重，普通人也不知道怎么用。"真正的 Open 是让技术触手可及。
\end{importantbox}

\subsection{本章小结}

OpenAI 面临的核心矛盾是\textbf{规模增长与创新效率}之间的张力。信息流通的一致性——无论是代码库、组织架构还是人际沟通——是翁家翌反复强调的解法。他甚至设想了一个极致方案：未来可能有一个拥有无限 context 的 AI agent 来当 CEO，负责所有的 context sharing 和 decision making。


